\documentclass[10pt,twoside]{article}
\usepackage[T1]{fontenc}
\usepackage[utf8]{inputenc}
\usepackage{amsmath,amssymb,amsthm,amsfonts,mathrsfs,mathtools,bm}
\usepackage{graphicx,xcolor,tikz,booktabs,multirow,tabularx,array,float}
\usepackage{caption,subcaption,enumitem,setspace,verbatim}
\usepackage{natbib,hyperref,fancyhdr,geometry,xspace,pdflscape}
\usepackage{algorithm,algpseudocode}
\geometry{papersize={16.5cm,23.5cm},top=2cm,bottom=1.5cm,left=1.5cm,right=1.5cm}
\hypersetup{colorlinks=true,linkcolor=blue,citecolor=magenta,urlcolor=cyan}
\pagestyle{fancy}
\fancyfoot{}
\renewcommand{\headrulewidth}{.01cm}
\fancyhead[RE,LO]{\thepage}
\newtheorem{dfn}{Definition}[section]
\newtheorem{thm}[dfn]{Theorem}
\newtheorem{pro}[dfn]{Proposition}
\newtheorem{rem}[dfn]{Remark}
\newtheorem{lem}[dfn]{Lemma}
\newtheorem{cor}[dfn]{Corollary}
\newtheorem{eg}[dfn]{Example}
\newtheorem{alg}{Algorithm}
\let\oldtitle\title
\renewcommand{\bibname}{References}
\numberwithin{equation}{section}
\graphicspath{{./}{main/p.1094/}{main/p.1094/Images/}{Images/}}
\makeatletter
\def\@maketitle{%
  \begin{center}
  \vspace*{1cm}
  {\Large\bfseries \@title\par}
  \vspace{0.8em}
  {\bfseries \@author\par}
  \vspace{0.5em}
  \end{center}
  \rule{\textwidth}{0.2mm}\\
}
\makeatother
\begin{document}
\thispagestyle{empty}

\begin{quote}


\includegraphics[width=5.98453in,height=0.99667in]{main//p.1094//vertopal_98ee39a136a74a2eadcf9766209bb0b0/media/image1.jpeg}

\end{quote}



\section{Method for Estimating Genotype Stability, Adaptability, and

Environmental Discrimination in Plant Breeding

Programs}



\textbf{Naser Sabaghnia}\footnote{\begin{quote}

  \textbf{Corresponding Author:} sabaghnia@maragheh.ac.ir

  \end{quote}}



\begin{quote}

\textbf{University of Maragheh, Iran}

\end{quote}



\textbf{Abstract:}



Breeding programs aim not only to improve crop performance but also to

ensure stability and adaptability across diverse environments.

Evaluating genotype resilience requires understanding both general

adaptability ability (GAA), the mean trait performance across multiple

environments, as well as specific adaptability ability (SAA), the

deviation of a genotype in a particular environment. A comprehensive

approach for assessing GAA, SAA, and genotype stability, is presented

for integrating both linear and nonlinear responses to environmental

variation. The method relies on multi-environment trials, partitioning

phenotypic variance into genotype, environment, and genotype ×

environment interaction components using ANOVA models, and calculating

stability indices such as the relative stability and stability value of

genotypes (SVG). This framework allows breeders to select genotypes with

optimal performance while maintaining stability, balancing trade-offs

between high mean trait values and environmental sensitivity.

Furthermore, it provides insights into environmental differentiation and

the compensatory or destabilizing effects of genotype × environment

interactions. By combining theoretical and practical aspects of adaptive

breeding, this approach offers a robust tool for guiding selection

strategies that enhance both productivity and resilience in crop

populations. To compare populations and select appropriate breeding

methods across different environments, a technique is introduced to

partition the phenotypic variance of a population into GAA and SAA

variances.



\textbf{Keywords:} Adaptability ability, Stability value, Genotype ×

environment interaction



\textbf{Mathematics Subject Classification (2010):} 62P10, 62J10, 62K10,

92D10.



\section{Introduction}



Plant breeders are challenged not only with the task of increasing plant

performance (productivity), but also with combining performance with

resistance to biotic and abiotic environmental conditions. Accordingly,

researchers have focused on studying and developing methods to evaluate

resilience (stability, general and specific adaptability of genotypes),

as well as breeding approaches aimed at optimizing these important plant

traits. Various approaches exist for estimating genotype stability. The

terms general and specific adaptability (Annicchiarico, 2002), which

reflect the overall response of a genotype across all environments and

its specific response in a particular environment, have been reported in

the literature. Studying genotypes across different environments can

also provide information about the environments themselves as a context

for breeding. Kilchevskii \& Khotilova, (1985), proposed a genetic

analysis method based on testing genotypes in multiple environments that

allows researchers to determine the general and specific adaptability

and stability of genotypes, and to compare environments in terms of

their ability to differentiate genotypes. In this approach, the breeding

value of a genotype can be estimated depending on the breeding

objectives, and the breeding process can be guided according to

adaptability levels. This method also enables the partitioning of a

population\textquotesingle s phenotypic variance into general and

specific adaptability variances, allowing comparisons among populations

and informed selection of breeding methods.



\section{General and specific adaptability

ability}



Adaptability is defined as the ability of a genotype to maintain its

specific phenotypic trait under a particular environmental condition.

General adaptability ability (GAA) of a genotype represents the mean

trait value across different environmental conditions, while specific

adaptability ability (SAA) represents the deviation from GAA in a

specific environment. The method for evaluating GAA and SAA is based on

testing a population of n genotypes across m environments. The number of

replications is k. Accordingly; \emph{x\textsubscript{ikr}} is the

phenotypic value of genotype i in replication r in environment k;

\emph{u} is the overall mean of all phenotypic values;

\emph{v\textsubscript{i}} is the effect of genotype i;

\emph{d\textsubscript{k}} is the effect of environment k;

\emph{(vd)\textsubscript{ik}} is the interaction effect of genotype i

with environment k; \emph{e\textsubscript{ikr}} is the random error

associated with phenotype ikr.



To determine the significance of genotype, environment, and genotype ×

environment interaction effects on the phenotypic variation of the

population, a two-way analysis of variance (ANOVA) is used.

Interpretation of the ANOVA results depends on the underlying

assumptions of the study. If genotypes and environments are non-random

samples, their effects are considered fixed, and information can be

obtained about specific genotypes and environments. In this case, Model

1 is used, in which all effect values, except for the independent

variable \emph{e\textsubscript{ikr}} are considered fixed. The variable

\emph{e\textsubscript{ikr}} is assumed to be normally distributed with

mean zero and variance σ\textsuperscript{2}\textsubscript{e}. The F-test

is used to determine the significance of genotype, environment, and

interaction effects. In this test, the corresponding mean squares are

compared to the mean square of the random deviations. If genotypes are a

random sample from the population and environments are a random sample

from all possible environments, then Model 2 is used. In this case, it

is possible to obtain information about the population parameters of

genotypes for a specific set of environments. In Model 2, only the

overall mean u is considered fixed, while the other effects are

independently and normally distributed with mean zero and variance

σ\textsuperscript{2}\textsubscript{θ}. In this case, to test hypotheses

regarding genotypes and environments, the mean square of genotype ×

environment interaction is used, and to estimate the genotype ×

environment interaction effect, the mean square of random deviations is

applied.



\section{Description of method}



Regarding accept the assumptions of Model 1, the effect of the general

adaptability ability (GAA) of genotype i is equal to u+v\emph{i}. The

deviation from the sum u+v\emph{i} represents the specific adaptability

ability effect (SAA\emph{\textsubscript{ik}}) of genotype i in

environment k. This effect consists of a linear component (the effect of

environment dk) and a nonlinear component (the interaction effect

(vd)\textsubscript{ik}). To determine the effects of

GAA\textsubscript{i} and SAA\textsubscript{ik}, the mean phenotypic

values across replications are calculated. In this way, the results of

testing n genotypes across m environments can be arranged in an n×m

table, where: x\textsubscript{i} is the sum of all phenotypic values for

genotype i, x\textsubscript{k} is the total phenotypic sum in

environment k, X.. is the grand total of all phenotypic values,

x\textsubscript{ik} is the phenotypic value of genotype i in environment

k. The GAA and SAA effects are calculated using the following formulas

(Fig. 1):



\[x_{i.} = \sum_{k}^{}x_{ik} = x_{i1} + x_{i2} + \ldots + x_{im}\]



\[x_{nk} = \sum_{k}^{}x_{ik} = x_{1k} + x_{2k} + \ldots + x_{nk}\]



\[x_{ik} = u + v_{i} + d_{k} + {(vd)}_{ik}\]



\[x_{ik} = u + {GAA}_{i} + {SAA}_{ik}\]



\begin{longtable}[]{@{}

  >{\raggedright\arraybackslash}p{(\columnwidth - 0\tabcolsep) * \real{1.0000}}@{}}

\toprule()

\begin{minipage}[b]{\linewidth}\raggedright

\includegraphics[width=4.23313in,height=1.83843in]{main//p.1094//vertopal_98ee39a136a74a2eadcf9766209bb0b0/media/image2.jpeg}

\end{minipage} \\

\midrule()

\endhead

Figure 1. Schematic illustration of stability evaluation \\

\bottomrule()

\end{longtable}



The variance of random deviations for the mean phenotypic value in a

specific environment x\textsubscript{ik} is computed as:



\[{var}_{xik} = \frac{\sigma_{e}^{2}}{c} = \sigma^{2}\]



The variance of the difference between the means of two genotypes is

determined as:



\[{var}_{(x_{ik}/x_{nm})} = 2\sigma^{2}\]



To obtain the standard error of the differences between effects, the

square root of the corresponding variance must be taken. The least

significant difference (LSD) is then obtained by multiplying the

standard error by a constant magnitude. Comparison of genotypes in terms

of general adaptability can be performed by comparing

GAA\textsubscript{i} values. To estimate the interaction strength of

genotype i with environments, the genotype × environment interaction

variance can be used. To measure the stability of genotype i, it is

recommended to use the variance of SAA\textsubscript{i}. Therefore,

stability can be defined as the capacity of a genotype to maintain its

phenotypic expression under various environmental conditions through

organized mechanisms.



\section{Stability analysis}



Stabilizing selection provides the evolutionary basis for developing

organized mechanisms that ensure the phenotypic stability of genotypes.

In this way, stability in the expression of one trait can be combined

with flexibility in the expression of other traits. The proposed method

for estimating stability is based on combining the linear and nonlinear

components of the genotype's response to the environment. This approach

differs from other established methods like regression model of Finlay

and Wilkinson (1963), in which stability is measured using a linear

response; ecovalence in which stability is estimated via a nonlinear

response; and regression models of Eberhart \& Russell (1966), in which

parameters for both linear and nonlinear genotype responses are

determined.



The concept of relative genotype stability is defined as follows:



\[s_{gi} = \frac{\sigma_{{GAA}_{i}}}{u + {SAA}_{i}} \times 100\]



This index allows comparison of experimental results obtained across

different sets of crops, genotypes, environments, and traits under

study. Essentially, relative stability is analogous to the coefficient

of variation when studying a genotype across multiple environments. The

ratio l\textsubscript{gi} can serve as a nonlinear index of genotype i's

response to the environment:



\[l_{gi} = \frac{\sigma_{GE}^{2}}{\sigma_{SAA}^{2}}\]



If lgi→1, the genotype exhibits mostly nonlinear responses across most

environments, while if lgi→0, the genotype's response to the environment

is predominantly linear. The response of a genotype to improved

environmental conditions can be quantified using the regression

coefficient of the genotype across environments. There is considerable

disagreement among breeders regarding the interpretation of genotype

stability. Consequently, Eberhart and Russell (1966), suggested that an

ideal genotype should have: a high mean trait value, a regression

coefficient close to 1, and the smallest deviation from the regression

line. However, the work of Mulder et al. (2007), demonstrated that

through breeding, any combination of mean trait value and environmental

sensitivity can be achieved. The most acceptable perspective is that of

Finlay and Wilkinson (1963), regression model which is defined an

optimal genotype as one that possesses, as high general adaptability,

maximum performance under favorable environments, and maximum stability.

According to Martínez-Peña et al. (2023), a genotype with moderate but

stable yield has higher economic value than a genotype with high

potential but highly variable yield. Similarly, Bailey-Serres et al.

(2019) reported that in crop production, only 30-50\% of the potential

productivity of dense genotypes is realized, primarily due to

insufficient environmental stability.



The final evaluation of a breeding genotype depends on the objectives of

adaptive breeding (Fig. 2). The main directions for adaptive breeding

are; breeding genotypes for specific adaptability (SAA) in a particular

environment. Such selection is only desirable under controlled and

managed environmental conditions (e.g., greenhouse, irrigated field). At

the same time, there is a risk of erosion of the genetic basis of

general viability and reduced genotype stability. Breeding genotypes for

general adaptability (GAA) across multiple environments, selection is

based on the mean phenotypic value across all environments. According to

Omel'yanyuk (2013), this approach achieves maximum trait enhancement

compared to selection in favorable or unfavorable environments, while

providing moderate environmental sensitivity. The latter is not

controlled during the breeding process.



\begin{longtable}[]{@{}

  >{\raggedright\arraybackslash}p{(\columnwidth - 0\tabcolsep) * \real{1.0000}}@{}}

\toprule()

\begin{minipage}[b]{\linewidth}\raggedright

\includegraphics[width=3.97546in,height=3.15064in]{main//p.1094//vertopal_98ee39a136a74a2eadcf9766209bb0b0/media/image3.jpeg}

\end{minipage} \\

\midrule()

\endhead

Figure 2. Flowchart of breeding and selection strategy \\

\bottomrule()

\end{longtable}



\section{Breeding for GAA while considering

stability}



The theoretical basis of this adaptive breeding approach is concerning

the inseparable forms of natural selection: directional and stabilizing.

The first form alters the structure and performance of organisms, while

the second stabilizes morphogenesis, leading to automated ontogeny and

the development of regulatory mechanisms. The effectiveness of

stabilizing selection in plants has been demonstrated in several studies

(Cresswell, 1998; Rolhauser \& Pucheta, 2017). However, stabilizing

selection alone cannot create new traits; it is only important for

preserving valuable population characteristics and increasing resistance

to environmental fluctuations. Breeding for GAA with consideration of

stability requires a specific selection criterion, which allows the

combination of performance and stability in the genotype.



Kilchevskii \& Khotilova, (1985), proposed the stability value genotype

(SVG) as:



\[{SVG}_{gi} = \frac{\sqrt{\left( \sum_{}^{}{({SAA}_{gi})}^{2} \right) - (\frac{MSE}{df})}}{u + {SAA}_{i}} \times 100\]



In practice, applying the usual selection criterion p=2 is unlikely,

because for large values of σ\textsubscript{SAAi}, negative SVG values

may occur, and some fertility genes could be lost. As a first step

toward achieving an optimal balance between performance and stability in

breeding, we suggest using the population mean SVG equal to 0.5 grand

mean. This facilitates establishing the relationship between the average

relative stability of genotypes and the value of p. In this approach, at

the lowest value, selection prioritizes stability, while at the highest

value, selection favors performance more intensely. Environmental

evaluation is also crucial in breeding. Environments can be classified

into three types; stable, in which population polymorphism does not

appear; analytical, which helps identify different biotypes; harmonized,

where the viability of different biotypes is constrained and differences

among them are equalized. The above analysis method allows not only the

evaluation of genotype adaptability parameters, but also defining the

environment as a breeding context.



When interpreting SAA and GAA variances, a question arises: how

informative are they, and can they replace genotype × environment

interaction variance? Some researchers believe that genotype ×

environment interaction variance can serve as a measure of genotype

stability. However, statistical evaluation does not reflect all

biological realities of genotype × environment interaction, including

increase or decrease in genotype (or environment) effects. If the signs

of genotype (environment) effects and the interaction do not match, the

variance of environmental differentiation or genotype stability

(compensation effect) decreases. If the signs are aligned, the

corresponding variance increases (destabilizing effect). Thus, even if

genotype × environment interaction variance is identical, the variance

of genotype stability and the variance of environmental differentiation

may differ. The practical importance of the compensation/destabilization

effect lies in the unpredictability of genotype stability or

environmental differentiation when using

σ\textsuperscript{2}\textsubscript{GE}, as a criterion for selecting

genotypes or environments. Therefore, SAA and GAA variances provide more

objective information on genotype stability and environmental

differentiation compared to genotype × environment interaction variance.



\section{Conclusion}



This study introduces an integrated framework for evaluating genotype

adaptability and stability based on general adaptability ability (GAA)

and specific adaptability ability (SAA). By partitioning phenotypic

variation into general and environment-specific components, the approach

provides a more informative assessment of genotype performance across

environments than conventional GE interaction analysis alone. The use of

GAA, SAA, and stability indices enables simultaneous selection for high

yield and stability, supporting more effective breeding decisions.



\section{References}



Annicchiarico, P. (2002). \emph{Genotype x environment interactions:

challenges and opportunities for plant breeding and cultivar

recommendations}. Food and Agriculture Organization of the United

Nations. Rome, 2002.



Bailey-Serres, J., Parker, J. E., Ainsworth, E. A., Oldroyd, G. E., \&

Schroeder, J. I. (2019). Genetic strategies for improving crop yields.

\emph{Nature,} 575(7781), 109-118.



Cresswell, J. E. (1998). Stabilizing selection and the structural

variability of flowers within species. \emph{Annals of Botany,} 81(4),

463-473.



Eberhart, S.A., \& Russell, W.A. (1966). Stability parameters for

comparing varieties. \emph{Crop Science,} 6(1), 36-40.



Finlay, K. W., \& Wilkinson, G. N. (1963). The analysis of adaptation in

a plant-breeding programme. \emph{Australian Journal of Agricultural

Research}, 14(6), 742-754.



Kilchevskiy, A.V., \& Khotyleva, L.V. (1985). Method of evaluation of

adaptive ability and stability of genotypes, the differentiating ability

of environment. \emph{Genetics,} 21(9), 1481-1490.



Martínez-Peña, R., Rezzouk, F. Z., del Carmen Díez-Fraile, M.,

Nieto-Taladriz, M. T., Araus, J. L., Aparicio, N., \& Vicente, R.

(2023). Genotype-by-environment interaction for grain yield and quality

traits in durum wheat: Identification of ideotypes adapted to the

Spanish region of Castile and León. \emph{European Journal of Agronomy,}

151, 126951.



Mulder, H. A., Bijma, P., \& Hill, W. G. (2007). Prediction of breeding

values and selection responses with genetic heterogeneity of

environmental variance. \emph{Genetics,} 175(4), 1895-1910.



Omel'yanyuk, L. V. (2013). Specificity of the response of pea varieties

of different morphotypes to changes in the hydrodynamic conditions

during the vegetation period. \emph{Russian Agricultural Sciences,}

39(3), 236-240.



Rolhauser, A. G., \& Pucheta, E. (2017). Directional, stabilizing, and

disruptive trait selection as alternative mechanisms for plant community

assembly. \emph{Ecology,} 98(3), 668-677.

\end{document}
